\section{From jets to the Seshadri constant}\label{sec:positivity}

The interpolation statement now gives the lower bound for the Seshadri
constant.  We first spell out the passage from jets to intersection numbers,
including at singular points of a curve, and then recall the matching volume
bound.  The relation between asymptotic jet separation and Seshadri constants
is classical; see \cite[Theorem~6.4]{Demailly1992} in the complex setting
and \cite[Section~2, Equation~(2.3)]{MustataSchwede2014} in arbitrary
characteristic.
The arguments below work in arbitrary characteristic.

\begin{lemma}[Jets and curve multiplicities]\label{lem:jet-curve-bound}
Let $Y$ be a smooth integral projective variety over an algebraically closed
field, let $A$ be a line bundle on $Y$, and let $q_1,\ldots,q_r$ be distinct
points.  Suppose that $H^0(Y,A)$ separates jets through order $m\geq 1$ at
these points.  Then every integral curve $C\subset Y$ meeting at least one
of them satisfies
\[
 A\cdot C\geq m\sum_{i=1}^r\mult_{q_i}C,
\]
where the multiplicity is zero at a point outside $C$.
\end{lemma}

\begin{proof}
Choose such a point $q_j$.  Write $R=\mathcal O_{C,q_j}$ and let $\mathfrak n$
be its maximal ideal.  The vector space
$\mathfrak n^m/\mathfrak n^{m+1}$ is nonzero: otherwise Nakayama's lemma would
give $\mathfrak n^m=0$, contrary to $R$ being a one-dimensional local domain.
In a local frame of $A$, lift a nonzero class in this vector space to an
ambient jet of order $m$ at $q_j$.  This lift can be chosen in the image of
$\mathfrak m_{Y,q_j}^m$.  Prescribe that jet at $q_j$ and the zero jet at all
other marked points.  Jet separation produces a section $\sigma\in H^0(Y,A)$
which vanishes to order at least $m$ at every $q_i$ and whose restriction to
$C$ is nonzero.

For completeness, the local multiplicity estimate used here is
\begin{equation}\label{eq:local-zero-length}
 \operatorname{length}_R(R/(f))\geq m\,e(R)
 \qquad(0\ne f\in\mathfrak n^m),
\end{equation}
where $e(R)$ denotes the Hilbert--Samuel multiplicity of a one-dimensional
Noetherian local domain $(R,\mathfrak n)$; see
\cite[Definition~43.15.1, Tag~0AZV]{StacksProject}.  Indeed, for positive integers $t$,
the nonzero divisor $f$ gives
\[
 t\operatorname{length}_R(R/(f))
 =\operatorname{length}_R(R/(f^t))
 \geq\operatorname{length}_R(R/\mathfrak n^{mt})
 =mt\,e(R)+O(1).
\]
Divide by $t$ and let $t$ tend to infinity.  Applying
\eqref{eq:local-zero-length} to the local equations of $\sigma|_C$ at the
marked points proves the inequality: the sum of all local zero lengths of a
nonzero section of $A|_C$ is $\deg(A|_C)=A\cdot C$.
\end{proof}

\begin{lemma}[Volume bound]\label{lem:volume-bound}
Let $Y$ be a smooth integral projective variety of dimension $d\geq2$ over an
algebraically closed field, let $A$ be ample, and let $q_1,\ldots,q_r$ be
distinct points.  Then
\[
 \inf_{C\cap\{q_1,\ldots,q_r\}\ne\varnothing}
 \frac{A\cdot C}{\sum_i\mult_{q_i}C}
 \leq \left(\frac{A^d}{r}\right)^{1/d},
\]
where the infimum is over integral curves in $Y$.
\end{lemma}

\begin{proof}
Denote the infimum by $\varepsilon$.  If $\varepsilon=0$, the conclusion
is immediate.  Otherwise, blow up the marked points:
\[
 \pi:\widetilde Y\longrightarrow Y,
 \qquad E_1,\ldots,E_r\subset\widetilde Y.
\]
For each rational number $0\leq s<\varepsilon$, the divisor class
\[
 D_s=\pi^*A-s\sum_i E_i
\]
is nef.  To check this on curves, first observe that
$D_s|_{E_i}=s\mathcal O_{\PP^{d-1}}(1)$, which has nonnegative degree on every
curve in $E_i$.  Every other integral curve in $\widetilde Y$ is the strict
transform $\widetilde C$ of an integral curve $C\subset Y$, and
\begin{equation}\label{eq:exceptional-multiplicity}
 E_i\cdot\widetilde C=\mult_{q_i}C.
\end{equation}
Thus $D_s\cdot\widetilde C\geq0$ by the definition of $\varepsilon$; for curves
missing all marked points this follows from the ampleness of $A$.

Here is a local justification of \eqref{eq:exceptional-multiplicity}, which
also covers singular curves.  If $q_i\in C$, the Rees algebra construction
identifies $\widetilde C$ with the blow-up of $C$ at the marked points.
Indeed, the ambient Rees algebra surjects onto the Rees algebra of the
restricted ideal.  The latter is integral, and its Proj agrees with $C$
away from the centers, so its closed image is exactly the strict transform.
For
$R=\mathcal O_{C,q_i}$ with maximal ideal $\mathfrak n$, its exceptional fiber
is
\[
 \operatorname{Proj}\bigl(\operatorname{gr}_{\mathfrak n}R\bigr).
\]
It is the effective Cartier divisor cut out on $\widetilde C$ by $E_i$.
Its length is $e(R)$: the Hilbert function
$\dim_{k(q_i)}(\mathfrak n^t/\mathfrak n^{t+1})$ equals $e(R)$ for all
sufficiently large $t$, and its eventual value is the length of this
zero-dimensional projective scheme.  This proves
\eqref{eq:exceptional-multiplicity}; if $q_i\notin C$, both sides vanish.

A nef class on a projective variety is a limit of ample classes and therefore
has nonnegative top self-intersection.  The exceptional divisors are
disjoint, $\pi^*A$ restricts trivially to each of them, and
$E_i^d=(-1)^{d-1}$.  Consequently
\[
 0\leq D_s^d=A^d-rs^d.
\]
If $\varepsilon>(A^d/r)^{1/d}$, even if $\varepsilon=+\infty$, there is a
rational number $s$ strictly between $(A^d/r)^{1/d}$ and $\varepsilon$.
The displayed inequality excludes such an $s$, proving the bound.
\end{proof}

\begin{proof}[Proof of Theorem~\ref{thm:main}]
Choose $r_0$ as in Proposition~\ref{prop:generic-jets} and fix $r\geq r_0$.
Put $V=L^n$ and $w=(V/r)^{1/n}$.  Work over $K_r$ at the geometric generic tuple
$p_1,\ldots,p_r$.  For every fixed real number $0<\theta<w$ and all
sufficiently large integers $N$, Proposition~\ref{prop:generic-jets} gives
separation of jets through order $m=\lfloor N\theta\rfloor$ by
$H^0(X_{K_r},L_{K_r}^{\otimes N})$ at this fixed tuple.

For large $N$ we have $m\geq1$, so Lemma~\ref{lem:jet-curve-bound} yields
\[
 \frac{L_{K_r}\cdot C}{\sum_i\mult_{p_i}C}
 \geq\frac{\lfloor N\theta\rfloor}{N}
\]
for every integral curve in the defining infimum.  Letting $N$ tend to
infinity gives $\varepsilon_r(X,L)\geq\theta$.  Since $\theta<w$ was
arbitrary, $\varepsilon_r(X,L)\geq w$.  Intersection numbers are unchanged
by extension of the algebraically closed ground field, so
$(L_{K_r})^n=V$; Lemma~\ref{lem:volume-bound} supplies the reverse inequality.
\end{proof}

\begin{remark}[Surfaces]\label{rem:surfaces}
The same proof gives the conclusion of Theorem~\ref{thm:main} when $n=2$.
Indeed, Propositions~\ref{prop:initial-simplex}, \ref{prop:reshape},
and~\ref{prop:generic-jets} all allow $n=2$; the reshaping is then performed
just once, with no unchanged coordinates.  Both intersection lemmas above
also apply in dimension two.
\end{remark}
